\section{信道编码与汉明码：让收到的 0/1 能自己纠错}

这一讲只干一件事：\textbf{让收到的 0/1 自己把错改回来}。学完你应该能拿一支笔，把 8 位信息
\texttt{10011100} 手工变成 12 位码字 \texttt{111100101100}；并且在路上被翻错一个 bit 之后，
只看收到的 12 位，就反推出错在第几位、把它翻回来。全程只需要你会三样东西：
数\textbf{数}、判\textbf{奇偶}、外加一个叫\textbf{异或}的运算，不涉及概率论。

\subsection{先讲个故事：为什么要"故意多写几位"}

室友在电话里让你帮他记一个 11 位手机号，线不太好，你听到的数字里可能有 1 位是错的。
他念完就挂了。问题来了：你手里这 11 位，\textbf{你完全不知道其中有没有错}。
因为手机号本来就"可以是"任何 11 位数字，你没有任何东西可以拿来核对。也就是说，
\textbf{线路确实可能翻错 1 个 bit，但收方根本发现不了}——这就是"直接发原码"的处境：
错误发生了，却没有留下任何痕迹；更别提确定是哪一位错了。

现实世界的答案是"多带几位没用的数字"，而且你天天在用：身份证号 18 位，前 17 位是信息，
\textbf{最后 1 位是校验码}，它不是新信息，只用来查前面 17 位有没有被抄错；书的 ISBN、
快递单号、银行卡号，最后一位通常也是这个角色。这些"多出来的位"有个正经名字：
\textbf{校验位}。它的取值不自由，而是\textbf{完全由信息位算出来}的，于是收发双方多了一条约定：
\emph{"你算出来的和我算出来的对不上，就说明路上出错了"}。这就把"发现错误"从不可能变成了可能。

本讲处理的就是教材 5.6 节那套办法：要把 8 位信息送过去，但路上每一位都可能以概率 $p$
被翻错（这个"会翻错的线"叫\textbf{二元对称信道}，简称 BSC；设 $p<1/2$，意思是线虽然会错，
但错得不算多——若 $p>1/2$，把收到的每一位都反过来反而更好，这种信道没有讨论价值）。
做法：\textbf{不直接发这 8 位，而是再补 4 位校验，一共发 12 位}。这 4 位不携带新信息，
纯粹是"为了纠错而花掉的开销"。

\begin{idea}{一句话记住：冗余 = 用更长的码换可靠}
\textbf{冗余}（redundancy）就是那些"重复的、能被猜出来"的部分。原码没有冗余，错一位无从判断；
多带 4 位校验，等于给了收方一把尺子，让它能量出"这串码字不对劲"。代价很直白：本来 8 位就够，
现在要占 12 个位置，线要多花 $12/8=1.5$ 倍时间。\textbf{可靠性是用长度换来的。}
\end{idea}

特别要注意：\textbf{信道编码和信源编码的目标正好相反}。

\begin{center}\small
\setlength{\tabcolsep}{5pt}
\begin{tabular}{@{}lll@{}}
\toprule
 & 信源编码 & 信道编码\\
\midrule
要解决的问题 & 数据太大，传得慢 & 数据会传错，不可靠\\
目标 & 让码\textbf{变短} & 让码\textbf{变长}\\
对待冗余 & 尽可能\textbf{挖掉} & 尽可能\textbf{加入}（而且要有规律）\\
衡量指标 & 平均码长、编码效率 & 码率 $R$、纠错能力 $t$\\
典型例子 & Huffman 编码 & 重复码、汉明码\\
\bottomrule
\end{tabular}
\end{center}

一句话：\textbf{信源编码是"把水挤干"，信道编码是"往水里掺沙子"}；掺的沙子必须有规律，
否则收方也不知道哪些是沙子。

\begin{defbox}{四个术语（后面全靠它们说话）}
\textbf{信息位} $D_1,D_2,\dots,D_k$：真正想传的内容，本讲 $k=8$；$D$ 是 data，
$D_1$ 即"第 1 个数据位"。\quad
\textbf{校验位} $C_1,C_2,\dots,C_r$：由信息位算出来的位，本讲 $r=4$；$C$ 是 check，
$C_1$ 即"第 1 个校验位"。\quad
\textbf{码长} $n=k+r$：一个码字共多少位，本讲 $n=8+4=12$。\quad
\textbf{码字}：这 $n$ 位合在一起的整体，比如 \texttt{111100101100} 就是一个 12 位码字。
\end{defbox}

\begin{qa}{先把符号扫盲一遍：$\oplus$ 是什么？}
叫\textbf{异或}。规则只有两条：两个数\textbf{相同}得 0、\textbf{不同}得 1，即
$0\oplus0=0$，$0\oplus1=1$，$1\oplus0=1$，$1\oplus1=0$。要用的就两个性质：
(1) \textbf{连加看奇偶}——一串 bit 异或起来得 1 表示"里面有奇数个 1"，得 0 表示"偶数个 1"；
(2) \textbf{自己消掉自己}——$x\oplus x=0$，$x\oplus 0=x$，所以 $x\oplus y=0$ 等价于 $x=y$。
这两条就是整套汉明码的地基。
\end{qa}

\subsection{三个必须先钉死的概念}

\subsubsection{码率 $R=k/n$}

\begin{defbox}{码率}
一次发 $n$ 位，其中只有 $k$ 位是真信息，则 $R=k/n$。
\end{defbox}

人话：\textbf{每占用 1 个信道位置，平均能送进去多少位真信息}。$R$ 越接近 1 越"划算"，
越小说明为了纠错花的代价越大。三个码比一比：

\begin{center}\small
\setlength{\tabcolsep}{6pt}
\begin{tabular}{@{}lccccl@{}}
\toprule
码 & $k$ & $r$ & $n$ & $R=k/n$ & 说明\\
\midrule
不编码（裸传） & 8 & 0 & 8 & $8/8=1$ & 一点不浪费，错一位就完蛋\\
3 重复码 & 1 & 2 & 3 & $1/3\approx0.333$ & 每位说三遍，最稳也最亏\\
本讲的汉明码 & 8 & 4 & 12 & $8/12=2/3\approx0.667$ & 用一半浪费换"纠 1 位错"\\
教材 (7,4) 汉明码 & 4 & 3 & 7 & $4/7\approx0.571$ & 教材 5.6.2 节的例子\\
\bottomrule
\end{tabular}
\end{center}

注意后两行：同样是"纠 1 位错"，汉明码的 $2/3$ 是重复码 $1/3$ 的两倍。
\textbf{汉明码值得学的理由不是它更可靠，而是它更划算。}

\subsubsection{汉明距离：两个码字"差几位"}

\begin{defbox}{汉明距离}
两个等长的 0/1 串 $x$、$y$ 逐位对齐比一比，\textbf{不同的位置有几个}，
就叫它们的\textbf{汉明距离}，记作 $d(x,y)$。就是"数不同"。
\end{defbox}

例：$x=1011$ 与 $y=1101$。

\begin{center}\small
\setlength{\tabcolsep}{6pt}
\begin{tabular}{@{}lcccc@{}}
\toprule
位置 & 第 1 位 & 第 2 位 & 第 3 位 & 第 4 位\\
\midrule
$x$ & 1 & 0 & 1 & 1\\
$y$ & 1 & 1 & 0 & 1\\
一样吗 & 一样 & \textbf{不一样} & \textbf{不一样} & 一样\\
\bottomrule
\end{tabular}
\end{center}

2 个位置不同，故 $d(1011,1101)=2$。三条必记的事实：$d(x,y)=0$ 当且仅当 $x=y$；
$d(x,y)=d(y,x)$；若发 $x$ 路上被翻了 $t$ 位收到 $y$，则 $d(x,y)=t$——
\textbf{"错了几位"和"汉明距离"是同一件事}。

对一个码（一堆码字组成的集合）还有个关键数字：\textbf{最小距离} $d_{\min}$，即把所有两两之间的
汉明距离都算一遍、取最小的那个。整个码的抗错能力由它决定。

\subsubsection{纠错能力：$d_{\min}\ge 2t+1$ 才能纠 $t$ 位错}

\begin{defbox}{纠错能力判据}
$d_{\min}\ge 2t+1$ 是"能保证纠正 $t$ 位错误"的充要条件。
反过来读最常用：$d_{\min}=3$ 能纠 1 位错，$d_{\min}=5$ 能纠 2 位，$d_{\min}=7$ 能纠 3 位。
\end{defbox}

为什么是 $2t+1$ 而不是 $t+1$？分三步看。

\textbf{第一步（场景）。} 发送码字 $x$，路上被翻 $t$ 位，收到 $y$，则 $d(x,y)=t$——
收方手里这个 $y$ 离真正的答案 $x$ 有 $t$ 步远。

\textbf{第二步（收方干什么）。} 收方不知道 $x$，只能拿 $y$ 去和码本里所有码字比，
挑\textbf{最近的}当答案（叫"最小距离译码"）。要成功，就得保证 $y$ 离 $x$ 比离任何别的码字都近。

\textbf{第三步（算）。} 设另一个码字 $x'$ 满足 $d(x,x')\ge d_{\min}$。三角不等式给出
$d(x,x')\le d(x,y)+d(y,x')$，于是
\[
d(y,x')\ \ge\ d(x,x')-d(x,y)\ \ge\ d_{\min}-t .
\]
要 $d(y,x')>d(x,y)=t$，即 $d_{\min}-t>t$，也就是 $d_{\min}>2t$；整数条件下写作
$d_{\min}\ge 2t+1$。

\textbf{人话版。} 这个 $2t+1$ 是把距离劈成三段用：

\begin{center}\small
\setlength{\tabcolsep}{6pt}
\begin{tabular}{@{}lll@{}}
\toprule
段 & 用掉多少 & 干什么\\
\midrule
给错误留的余地 & $t$ & 收到的东西可能已被踹出去 $t$ 步，这段是"宽容区"\\
给对手留的余地 & $t$ & 别的码字也得允许别人误认到 $t$ 步以内，不能挤着\\
必须多出的 & $1$ & 保证"我最先赢"是\textbf{严格}赢，不是打平\\
\bottomrule
\end{tabular}
\end{center}

若只留 $2t$ 步，就会出现"$x$ 错了 $t$ 位、$x'$ 也错了 $t$ 位，两边正好撞在同一个 $y$ 上"，
收方彻底没主意、只能瞎猜——\textbf{所以必须再拉开 1 步}，这就是"留 2 倍余量再多 1"的来历。
$d_{\min}=3$ 时 $2t+1=3$ 给出 $t=1$：错 1 位能纠，错 2 位不敢保证（可能离另一个码字更近，
纠错反而纠歪）。汉明码"只纠 1 位错"不是它不努力，是它只有 3 的距离。

\subsection{重复码：最笨但最管用的办法}

\textbf{编码规则。} 每个 bit 说三遍：想发 \texttt{0} 就发 \texttt{000}，想发 \texttt{1} 就发
\texttt{111}。码本只有两个码字 $\{$\texttt{000},\texttt{111}$\}$，$n=3$、$k=1$、$R=1/3$。
两者汉明距离是 3（三位全不同），由 $3\ge 2\times1+1$ 知\textbf{能纠 1 位错}。

\textbf{译码规则：多数表决}——收到 3 位，\textbf{哪个数字多就判哪个}：

\begin{center}\small
\setlength{\tabcolsep}{6pt}
\begin{tabular}{@{}cl@{}}
\toprule
收到的 3 位 & 判成\\
\midrule
\texttt{000} & \texttt{0}（没错）\\
\texttt{001}、\texttt{010}、\texttt{100} & \texttt{0}（纠了 1 位错）\\
\texttt{011}、\texttt{101}、\texttt{110} & \texttt{1}（纠了 1 位错）\\
\texttt{111} & \texttt{1}（没错）\\
\bottomrule
\end{tabular}
\end{center}

看第二行：本来发 \texttt{000}，第 3 位被翻成 \texttt{001}，收方数出"0 有两个、1 只有一个"，
判 \texttt{0}，\textbf{纠错成功}。这就是多数表决的全部内容。

\begin{calc}{算重复码的误码率（代入 $p=0.01$，算术一步不跳）}
设每位独立地以概率 $p$ 被翻错。

\textbf{（1）什么时候表决会出错？} 只有一种情况：3 位里\textbf{至少错 2 位}
（错的位数不超过 1，正确的数字一定占多数）。

\textbf{（2）分类计数。}"恰好错 2 位"：从 3 个位置里挑 2 个出错，有 $\binom{3}{2}=3$ 种挑法，
每种概率是"这 2 位错（$p\cdot p$）、剩下 1 位对（$1-p$）"，即 $p^{2}(1-p)$，
合计 $3p^{2}(1-p)$。"3 位全错"只有 1 种情况，概率 $p^{3}$。两类互斥，相加：
\[
P_e=3p^{2}(1-p)+p^{3}.
\]
人话：\textbf{误码率 = "错俩"的概率 + "错仨"的概率}。

\textbf{（3）化简}（考试写出这步会加分）：$P_e=3p^{2}-3p^{3}+p^{3}=3p^{2}-2p^{3}$。

\textbf{（4）代入 $p=0.01$。} 先算中间量：
\[
p^{2}=0.01\times0.01=10^{-4},\qquad p^{3}=p^{2}\cdot p=10^{-4}\times10^{-2}=10^{-6}.
\]
第一项 $3p^{2}(1-p)=3\times10^{-4}\times0.99=2.97\times10^{-4}$；
第二项 $p^{3}=10^{-6}=0.01\times10^{-4}$。相加：
\[
P_e=2.97\times10^{-4}+0.01\times10^{-4}=2.98\times10^{-4}.
\]

\textbf{（5）对比原来。} $P_e/p=2.98\times10^{-4}\big/10^{-2}=2.98\times10^{-2}\approx3\%$，
即错误概率降到原来的约\textbf{三十分之一}。

\textbf{（6）代价。} 码率从 1 掉到 $1/3$：传同样多的信息要多花 3 倍时间。推广到"每 bit 发
$2m+1$ 遍"，则 $R=1/(2m+1)$，误码率趋于 0 但码率也趋于 0——\textbf{这条路走不通}。
\end{calc}

\begin{idea}{重复码的教训}
重复码靠"把信息抄 3 遍"，思路对但极浪费。汉明码的高明之处在于：\textbf{它不抄信息，
而是让 $r$ 个校验位各自"监视一片位置"，再把 $r$ 个监视结果按二进制拼起来，
拼出的数字直接就是"错在第几位"}。于是 $r$ 位校验能照看约 $2^{r}$ 个位置，而不是 1 个。
\end{idea}

\subsection{汉明码的构造（本讲核心，四步模板）}

下面这套流程照着抄，任何"$k$ 位信息、偶校验、构造汉明码"的题都能做。
例题即作业那道：\textbf{信息位 $k=8$，采用偶校验}。先把偶校验讲清楚：

\begin{defbox}{偶校验 / 奇校验}
\textbf{偶校验}：要求"一组 bit 里 1 的个数是\textbf{偶数}"，用异或写就是"这组 bit 异或起来
等于 0"。\textbf{奇校验}：要求"1 的个数是\textbf{奇数}"，即异或起来等于 1。
两者只是目标奇偶不同，其余全部照搬；本讲若不特别说明一律用偶校验。
\end{defbox}

\subsubsection{第一步：定校验位个数 $r$}

$r$ 不能随便挑，要满足
\[
2^{r}\ \ge\ n+1,\qquad\text{其中 } n=k+r .
\]
$n$ 里含 $r$，是自己包含自己的式子，标准做法是\textbf{逐个试}。

\textbf{人话翻译。} $r$ 个校验位最后给出 $r$ 个 0/1 答案，拼起来是 $r$ 位二进制，最多表示
$2^{r}$ 种情况；而收方要区分的状况有 $n+1$ 种：\textbf{"没错"}（1 种）加
\textbf{"第 1 位错"}、$\dots$、\textbf{"第 $n$ 位错"}（$n$ 种）。要能一一报到，就得
$2^{r}\ge n+1$。把 $n=8+r$ 代进去逐个试：

\begin{center}\small
\setlength{\tabcolsep}{6pt}
\begin{tabular}{@{}ccccccl@{}}
\toprule
试的 $r$ & $n=8+r$ & 需区分 $n+1$ 种 & 可用状态 $2^{r}$ & 够不够\\
\midrule
$r=1$ & $n=9$ & $10$ & $2^{1}=2$ & $2<10$，不够\\
$r=2$ & $n=10$ & $11$ & $2^{2}=4$ & $4<11$，不够\\
$r=3$ & $n=11$ & $12$ & $2^{3}=8$ & $8<12$，不够\\
$r=4$ & $n=12$ & $13$ & $2^{4}=16$ & $16\ge13$，\textbf{可行}\\
\bottomrule
\end{tabular}
\end{center}

故 \fbox{$r=4$}，\fbox{$n=k+r=8+4=12$}，这是一个 \textbf{(12,8) 码}，码率 $R=8/12=2/3$。

交叉验证一下教材：$k=4$ 时 $r=2$ 得 $n=6$、$2^{2}=4<7$ 不够；$r=3$ 得 $n=7$、$2^{3}=8\ge8$
恰好够，故 (7,4) 汉明码，与教材 "$n=2^{r}-1$、$k=2^{r}-1-r$" 一致。我们这题 $k=8$ 不是这种
"刚好填满"的值，所以 $2^{4}=16$ 有富余（只用了 13 个状态），13、14、15 号位置不存在，
永远不会被报到，不影响使用。

\subsubsection{第二步：定位置（校验位放在 1、2、4、8 位）}

\begin{defbox}{位置规则（必须背下来）}
把码字的 $n$ 个位置从 1 开始编号。\textbf{编号是 2 的整数次幂的位置放校验位}，即第
$1,2,4,8,16,\dots$ 位（也就是 $2^{0},2^{1},2^{2},2^{3},2^{4},\dots$）放 $C_1,C_2,C_4,C_8,\dots$；
\textbf{其余位置按从小到大的顺序}放信息位 $D_1,D_2,\dots,D_k$。
\end{defbox}

注意校验位的\textbf{下标就是它所在的位置号}：$C_1$ 在第 1 位，$C_4$ 在第 4 位，$C_8$ 在第 8 位。
好处是"名字即地址"，纠错时不容易搞错。把 $n=12$ 排出来：

\begin{center}\small
\setlength{\tabcolsep}{4pt}
\begin{tabular}{@{}l*{12}{c}@{}}
\toprule
位置（十进制） & 1&2&3&4&5&6&7&8&9&10&11&12\\
\midrule
位置（4 位二进制地址） & 0001&0010&0011&0100&0101&0110&0111&1000&1001&1010&1011&1100\\
\midrule
这一位是谁 & $C_1$&$C_2$&$D_1$&$C_4$&$D_2$&$D_3$&$D_4$&$C_8$&$D_5$&$D_6$&$D_7$&$D_8$\\
\bottomrule
\end{tabular}
\end{center}

\textbf{读法。} 第 1、2、4、8 位是校验位 $C_1,C_2,C_4,C_8$（下标即位置号）；剩下的第
3、5、6、7、9、10、11、12 位是信息位，依次是 $D_1$ 到 $D_8$。

\textbf{这里的坑，老师说第 7 位和第 10 位最常写错}：因为第 8 位被 $C_8$ 占掉了，
所以 \textbf{$D_5$ 不是第 8 位，而是第 9 位}，往后全部顺延一位。一旦从这里错位，
$D_5\sim D_8$ 的位置会全错，整个码字就废了。口诀：\textbf{"逢 2 的幂次让给校验位，
数据位见空就填"}。

\subsubsection{第三步：写校验方程（按二进制地址分组）}

\textbf{分组规矩。} 每个位置号写成 4 位二进制（因为 $r=4$）。地址的\textbf{第 1 位}（从右数，
权值 1）是 1 的进 $G_1$、由 $C_1$ 管；\textbf{第 2 位}（权值 2）是 1 的进 $G_2$、由 $C_2$ 管；
\textbf{第 3 位}（权值 4）是 1 的进 $G_4$、由 $C_4$ 管；\textbf{第 4 位}（权值 8）是 1 的进
$G_8$、由 $C_8$ 管。照位置表逐位检查，得四组：

\begin{center}\small
\setlength{\tabcolsep}{4pt}
\begin{tabular}{@{}lll@{}}
\toprule
组 & 包含的位置 & 成员\\
\midrule
$G_1$（地址第 1 位为 1） & $\{1,3,5,7,9,11\}$ & $C_1,\ D_1,\ D_2,\ D_4,\ D_5,\ D_7$\\
$G_2$（地址第 2 位为 1） & $\{2,3,6,7,10,11\}$ & $C_2,\ D_1,\ D_3,\ D_4,\ D_6,\ D_7$\\
$G_4$（地址第 3 位为 1） & $\{4,5,6,7,12\}$ & $C_4,\ D_2,\ D_3,\ D_4,\ D_8$\\
$G_8$（地址第 4 位为 1） & $\{8,9,10,11,12\}$ & $C_8,\ D_5,\ D_6,\ D_7,\ D_8$\\
\bottomrule
\end{tabular}
\end{center}

\textbf{动手核对，别信脑袋。} 拿 $G_2$ 说：$2=0010$、$3=0011$、$6=0110$、$7=0111$、
$10=1010$、$11=1011$，第 2 位都是 1，全部进组；而 $1=0001$、$4=0100$、$5=0101$、
$8=1000$、$9=1001$、$12=1100$ 的第 2 位都是 0，全部不进。一个不多一个不少。

\textbf{偶校验方程}：每组里 1 的个数必须是偶数，即每组异或起来等于 0：
\begin{align*}
G_1:\quad & C_1\oplus D_1\oplus D_2\oplus D_4\oplus D_5\oplus D_7=0,\\
G_2:\quad & C_2\oplus D_1\oplus D_3\oplus D_4\oplus D_6\oplus D_7=0,\\
G_4:\quad & C_4\oplus D_2\oplus D_3\oplus D_4\oplus D_8=0,\\
G_8:\quad & C_8\oplus D_5\oplus D_6\oplus D_7\oplus D_8=0 .
\end{align*}
人话：\textbf{每一行是一个"检查组"，组里的位凑在一起，1 的个数必须凑成偶数}；
校验位的任务就是"补足那个差额"——组里信息位已有奇数个 1 则 $C$ 取 1，
已有偶数个 1 则 $C$ 取 0。由 $x\oplus y=0\Leftrightarrow x=y$，把 $C$ 解出来：
\begin{align*}
C_1&=D_1\oplus D_2\oplus D_4\oplus D_5\oplus D_7,\\
C_2&=D_1\oplus D_3\oplus D_4\oplus D_6\oplus D_7,\\
C_4&=D_2\oplus D_3\oplus D_4\oplus D_8,\\
C_8&=D_5\oplus D_6\oplus D_7\oplus D_8 .
\end{align*}
（在异或里"移项"就是两边同时异或同一个东西，所以把 $C$ 挪到等号右边即可。）

\subsubsection{第四步：代入信息码 \texttt{10011100}，求四个校验位}

题目给的 8 位信息码 \texttt{10011100}（最高位是 $D_1$），逐位拆开：

\begin{center}\small
\setlength{\tabcolsep}{6pt}
\begin{tabular}{@{}lcccccccc@{}}
\toprule
符号 & $D_1$ & $D_2$ & $D_3$ & $D_4$ & $D_5$ & $D_6$ & $D_7$ & $D_8$\\
\midrule
取值 & 1 & 0 & 0 & 1 & 1 & 1 & 0 & 0\\
\bottomrule
\end{tabular}
\end{center}

\begin{calc}{手把手求 $C_1,C_2,C_4,C_8$（异或中间结果全写出）}
\textbf{求 $C_1$}（$G_1$ 组，$1,3,5,7,9,11$ 位）：
\[
C_1=D_1\oplus D_2\oplus D_4\oplus D_5\oplus D_7=1\oplus0\oplus1\oplus1\oplus0 .
\]
从左往右一个个吃掉：$1\oplus0=1$；$1\oplus1=0$；$0\oplus1=1$；$1\oplus0=1$。
故 $C_1=1$。

\textbf{求 $C_2$}（$G_2$ 组，$2,3,6,7,10,11$ 位）：
\[
C_2=D_1\oplus D_3\oplus D_4\oplus D_6\oplus D_7=1\oplus0\oplus1\oplus1\oplus0 .
\]
和上面碰巧是同一串数，同样五步得 $C_2=1$。（组里 $D_1,D_4,D_6$ 三个 1，奇数个，
所以 $C_2$ 取 1 凑偶数。）

\textbf{求 $C_4$}（$G_4$ 组，$4,5,6,7,12$ 位）：
\[
C_4=D_2\oplus D_3\oplus D_4\oplus D_8=0\oplus0\oplus1\oplus0 .
\]
$0\oplus0=0$；$0\oplus1=1$；$1\oplus0=1$。故 $C_4=1$。（组里只有 $D_4$ 一个 1。）

\textbf{求 $C_8$}（$G_8$ 组，$8,9,10,11,12$ 位）：
\[
C_8=D_5\oplus D_6\oplus D_7\oplus D_8=1\oplus1\oplus0\oplus0 .
\]
$1\oplus1=0$；$0\oplus0=0$；$0\oplus0=0$。故 $C_8=0$。
（组里已有 $D_5,D_6$ 两个 1，偶数个，不用再添。）

\textbf{汇总}：\fbox{$C_1=1$}\quad\fbox{$C_2=1$}\quad\fbox{$C_4=1$}\quad\fbox{$C_8=0$}
\end{calc}

\textbf{拼成 12 位码字}：按位置表一位一位填格。

\begin{center}\small
\setlength{\tabcolsep}{4pt}
\begin{tabular}{@{}l*{12}{c}@{}}
\toprule
位置 & 1&2&3&4&5&6&7&8&9&10&11&12\\
\midrule
身份 & $C_1$&$C_2$&$D_1$&$C_4$&$D_2$&$D_3$&$D_4$&$C_8$&$D_5$&$D_6$&$D_7$&$D_8$\\
取值 & 1&1&1&1&0&0&1&0&1&1&0&0\\
\bottomrule
\end{tabular}
\end{center}

从左到右读出来：$\boxed{\ \texttt{1111\,0010\,1100}\ }$，即 \fbox{\texttt{111100101100}}。

\begin{calc}{回代验算：四组里 1 的个数是否都为偶数？（不写这步不给分）}
把校验位\textbf{也算进去}，数每组 1 的个数。

$G_1=\{1,3,5,7,9,11\}$，取值 $C_1,D_1,D_2,D_4,D_5,D_7=1,1,0,1,1,0$，
1 的个数 $=1+1+0+1+1+0=4$，\textbf{偶数} $\checkmark$

$G_2=\{2,3,6,7,10,11\}$，取值 $C_2,D_1,D_3,D_4,D_6,D_7=1,1,0,1,1,0$，
1 的个数 $=4$，\textbf{偶数} $\checkmark$

$G_4=\{4,5,6,7,12\}$，取值 $C_4,D_2,D_3,D_4,D_8=1,0,0,1,0$，
1 的个数 $=1+0+0+1+0=2$，\textbf{偶数} $\checkmark$

$G_8=\{8,9,10,11,12\}$，取值 $C_8,D_5,D_6,D_7,D_8=0,1,1,0,0$，
1 的个数 $=0+1+1+0+0=2$，\textbf{偶数} $\checkmark$

四组分别是 $4/4/2/2$，全为偶数，验算通过。
\end{calc}

\subsubsection{反过来用：收到的东西错在哪一位？}

设发送方发出 \texttt{111100101100}，路上第 10 位被咬了一口，$1$ 翻成 $0$，收方拿到
\[
y=\texttt{1111\,0000\,1100}\quad(\text{第 10 位由 }1\text{ 变成 }0).
\]

\begin{calc}{手算反推出错位置（综合征 / 伴随式）}
收方\textbf{并不知道}原码字，只有 $y$ 和那四条校验方程。做法：拿 $y$ 重新算四个组的奇偶性，
\textbf{算出来是奇记 1、是偶记 0}，四个 0/1 合起来叫\textbf{综合征}（伴随式），记
$(S_1,S_2,S_4,S_8)$。$y$ 从左到右十二位是 $1,1,1,1,0,0,1,0,1,0,0,0$。

\textbf{$S_1$（看 $G_1=\{1,3,5,7,9,11\}$）}：
$1\oplus1\oplus0\oplus1\oplus1\oplus0$，即 $1\oplus1=0$，$0\oplus0=0$，
$0\oplus1=1$，$1\oplus1=0$，$0\oplus0=0$，得 $S_1=0$（组内 4 个 1，偶数，没被破坏）。

\textbf{$S_2$（看 $G_2=\{2,3,6,7,10,11\}$）}：
$1\oplus1\oplus0\oplus1\oplus0\oplus0$，即 $1\oplus1=0$，$0\oplus0=0$，
$0\oplus1=1$，$1\oplus0=1$，$1\oplus0=1$，得 $S_2=1$（组内 3 个 1，奇数，\textbf{被破坏}）。

\textbf{$S_4$（看 $G_4=\{4,5,6,7,12\}$）}：
$1\oplus0\oplus0\oplus1\oplus0$，即 $1\oplus0=1$，$1\oplus0=1$，$1\oplus1=0$，
$0\oplus0=0$，得 $S_4=0$（组内 2 个 1，偶数，没被破坏）。

\textbf{$S_8$（看 $G_8=\{8,9,10,11,12\}$）}：
$0\oplus1\oplus0\oplus0\oplus0$，即 $0\oplus1=1$，$1\oplus0=1$，$1\oplus0=1$，
$1\oplus0=1$，得 $S_8=1$（组内 1 个 1，奇数，\textbf{被破坏}）。

\textbf{拼地址}：按权值 $1,2,4,8$ 把四个结果装成十进制位置号，
\[
S_1\cdot1+S_2\cdot2+S_4\cdot4+S_8\cdot8=0\times1+1\times2+0\times4+1\times8=10,
\]
写成二进制即 $(S_8S_4S_2S_1)=(1010)_{2}=10$。

\textbf{结论：第 10 位错了。} 把 $y_{10}$ 由 0 翻回 1，得 \texttt{111100101100}，
与原码字一模一样。\textbf{纠错完成，收方不需要回问发送方。}
\end{calc}

\begin{defbox}{综合征为什么能"报门牌号"}
因为\textbf{校验矩阵的每一列就是该位置的二进制地址}（见后面 $H$ 那一节）。某一位出错，
只会让"地址里为 1 的那几个组"的奇偶性翻掉，于是综合征正好等于那一列的地址。
所以\textbf{综合征 = 出错位置的二进制地址，直接读出来就行}，不用查表、不用试。
\end{defbox}

两种边界情形要会区分：综合征 $(0,0,0,0)$ 说明\textbf{没有错}（或错了至少 2 位恰好互相抵消，
汉明码管不了）；综合征指向第 1、2、4、8 位（地址 $0001,0010,0100,1000$）说明错的是
\textbf{校验位自己}，信息位没坏，把那个校验位翻回来即可，或干脆忽略。

\subsubsection{对照练习：老师课件里的那道例题}

课件与作业讲评那道题是：8 位信息码 \texttt{11101100}（最高位为 $D_1$），偶校验，构造汉明码。
做法\emph{一个字都不用改}，只换数据。数据位 $D_1$ 到 $D_8$ 依次取 $1,1,1,0,1,1,0,0$，
代入第三步解出的四式：
\begin{align*}
C_1&=D_1\oplus D_2\oplus D_4\oplus D_5\oplus D_7=1\oplus1\oplus0\oplus1\oplus0=1,\\
C_2&=D_1\oplus D_3\oplus D_4\oplus D_6\oplus D_7=1\oplus1\oplus0\oplus1\oplus0=1,\\
C_4&=D_2\oplus D_3\oplus D_4\oplus D_8=1\oplus1\oplus0\oplus0=0,\\
C_8&=D_5\oplus D_6\oplus D_7\oplus D_8=1\oplus1\oplus0\oplus0=0 .
\end{align*}
（前两行：$1\oplus1=0$，$0\oplus0=0$，$0\oplus1=1$，$1\oplus0=1$，得 1；后两行：
$1\oplus1=0$，$0\oplus0=0$，$0\oplus0=0$，得 0。）故 $C_1=1,C_2=1,C_4=0,C_8=0$，填格：

\begin{center}\small
\setlength{\tabcolsep}{4pt}
\begin{tabular}{@{}l*{12}{c}@{}}
\toprule
位置 & 1&2&3&4&5&6&7&8&9&10&11&12\\
\midrule
身份 & $C_1$&$C_2$&$D_1$&$C_4$&$D_2$&$D_3$&$D_4$&$C_8$&$D_5$&$D_6$&$D_7$&$D_8$\\
取值 & 1&1&1&0&1&1&0&0&1&1&0&0\\
\bottomrule
\end{tabular}
\end{center}

得 \fbox{\texttt{111011001100}}，与课件答案一致。回代验算（逐组数 1 的个数）：
\begin{itemize}[leftmargin=1.6em,itemsep=0pt]
  \item $G_1$：$C_1,D_1,D_2,D_4,D_5,D_7$ 取 $1,1,1,0,1,0$，共 4 个 1，偶数；
  \item $G_2$：$C_2,D_1,D_3,D_4,D_6,D_7$ 取 $1,1,1,0,1,0$，共 4 个 1，偶数；
  \item $G_4$：$C_4,D_2,D_3,D_4,D_8$ 取 $0,1,1,0,0$，共 2 个 1，偶数；
  \item $G_8$：$C_8,D_5,D_6,D_7,D_8$ 取 $0,1,1,0,0$，共 2 个 1，偶数。
\end{itemize}
四组全偶，验算通过。

\textbf{讲评提醒}：老师说这道题"汉明码第 7 位和第 10 位编码错误"是常见失分点。对照上表，
第 7 位是 $D_4=0$、第 10 位是 $D_6=1$。根源就是第二步那个坑：第 8 位被 $C_8$ 占了，
$D_5\sim D_8$ 整体后移一位。\textbf{填格之前先把位置表抄一遍、填完再对一遍，就能避开。}

\subsection{线性分组码的一点皮毛：$G$ 和 $H$}

上面的构造看似手工，背后是线性代数。记住两个矩阵的名字和用法就够用了。

\subsubsection{校验矩阵 $H$：把"每组异或为 0"写成一个矩阵}

把第三步的四个组直接写成矩阵：行是组、列是位置，属于该组的位置填 1、不属于填 0，
就得到\textbf{校验矩阵 $H$}（$4\times12$；这样的宽矩阵用表来排更清楚，数的是同一个东西）：

\begin{center}\small
\setlength{\tabcolsep}{4pt}
\begin{tabular}{@{}c*{12}{c}@{}}
\toprule
 & 1&2&3&4&5&6&7&8&9&10&11&12\\
\midrule
$G_1$ 行 & 1&0&1&0&1&0&1&0&1&0&1&0\\
$G_2$ 行 & 0&1&1&0&0&1&1&0&0&1&1&0\\
$G_4$ 行 & 0&0&0&1&1&1&1&0&0&0&0&1\\
$G_8$ 行 & 0&0&0&0&0&0&0&1&1&1&1&1\\
\bottomrule
\end{tabular}
\end{center}

读法：第 1 行对应 $G_1$，它哪几列是 1 就说明 $G_1$ 管哪几位（第 $1,3,5,7,9,11$ 列），
和前面完全一致——一行顶替了一句话。

\textbf{最重要的观察}：竖着看，\textbf{第 $i$ 列就是 $i$ 的 4 位二进制地址}（自下而上是高位
到低位）。如第 10 列是 $(0,1,0,1)^{\mathsf T}$，正好是 $10=8+2$；第 12 列是
$(0,0,1,1)^{\mathsf T}$，正好是 $12=8+4$。\textbf{$H$ 的第 $i$ 列 $=i$ 的二进制表示}——
这就是汉明码的全部秘密。有了 $H$，"12 位向量 $c$ 是合法码字"有了判据：
\[
cH^{\mathsf T}=\mathbf{0}\quad(\text{即四个组异或全为 }0).
\]
人话：\textbf{码字乘校验矩阵的转置得全零，就说明四条方程全满足}。
（所有加法乘法都在二元域上做：$1+1=0$、$1\times1=1$，也就是"异或"和"与"。）

\subsubsection{生成矩阵 $G$：把"信息位 $\to$ 码字"写成一个乘法}

$G$ 是 $8\times12$ 矩阵，\textbf{每一行是一个"基底码字"}：第 $i$ 行是"只有 $D_i=1$、
其余信息位为 0"时算出来的完整 12 位码字。有了这 8 行，任何信息向量都能拼出来：
\[
c=d\,G,\qquad d=(D_1,D_2,\dots,D_8)\in\{0,1\}^{8}.
\]
人话：\textbf{码字 = 信息向量乘生成矩阵}；$d$ 里哪几位是 1，就把 $G$ 对应的几行加起来（异或）。

\begin{center}\small
\setlength{\tabcolsep}{4pt}
\begin{tabular}{@{}c*{12}{c}@{}}
\toprule
 & 1&2&3&4&5&6&7&8&9&10&11&12\\
\midrule
$D_1$ 行 & 1&1&1&0&0&0&0&0&0&0&0&0\\
$D_2$ 行 & 1&0&0&1&1&0&0&0&0&0&0&0\\
$D_3$ 行 & 0&1&0&1&0&1&0&0&0&0&0&0\\
$D_4$ 行 & 1&1&0&1&0&0&1&0&0&0&0&0\\
$D_5$ 行 & 1&0&0&0&0&0&0&1&1&0&0&0\\
$D_6$ 行 & 0&1&0&0&0&0&0&1&0&1&0&0\\
$D_7$ 行 & 1&1&0&0&0&0&0&1&0&0&1&0\\
$D_8$ 行 & 0&0&0&1&0&0&0&1&0&0&0&1\\
\bottomrule
\end{tabular}
\end{center}

\textbf{用法示范（顺便复核前面的答案）。} $d=\texttt{10011100}$ 即 $D_1,D_4,D_5,D_6$ 为 1，
把这四行竖着异或：第 1 位 $1\oplus1\oplus1\oplus0=1$、第 2 位 $1\oplus1\oplus0\oplus1=1$、
第 3 位 $1\oplus0\oplus0\oplus0=1$、第 4 位 $0\oplus1\oplus0\oplus0=1$、
第 7 位 $0\oplus1\oplus0\oplus0=1$、第 9 位 $0\oplus0\oplus1\oplus0=1$、
第 10 位 $0\oplus0\oplus0\oplus1=1$；其余位置这四行都是 0，故异或为 0。
结果 $C_1=1,C_2=1,C_4=1,C_8=0$，码字仍为 \texttt{111100101100}，
\textbf{与第四步手算完全一致}。

\textbf{两个矩阵的关系。} $GH^{\mathsf T}=\mathbf{0}$。人话：\textbf{$G$ 的每一行都是合法码字，
而合法码字过一遍 $H$ 一定是全零}，所以乘出来是零矩阵。这条式子就是"生成矩阵"与
"校验矩阵"的定义关系；考试让你验证时逐行验（每行与 $H$ 各行的内积都为 0）即可。

\subsubsection{用 $H$ 算综合征：$S=yH^{\mathsf T}$}

收到 $y$ 后算 $S=yH^{\mathsf T}=(S_1,S_2,S_4,S_8)$。人话：\textbf{就是把 $y$ 代进四个校验
方程，把四个奇偶结果排成一个 4 位向量}。设错误图样 $e$ 只在第 $i$ 位为 1（其余为 0），
则 $y=c\oplus e$，于是
\[
S=yH^{\mathsf T}=(c\oplus e)H^{\mathsf T}
=\underbrace{cH^{\mathsf T}}_{=\,\mathbf{0}}\oplus\,eH^{\mathsf T}
=H\text{ 的第 }i\text{ 列}.
\]
人话：\textbf{综合征只跟"错在哪"有关，跟"发了什么"完全无关}——所以同一位置的错，
不管发的是什么码字，综合征都一样，这就是它能当"门牌号"的原因。于是 $S=\mathbf{0}$ 表示
无错；$S\ne\mathbf{0}$ 时 $S$ 就是 $H$ 的某一列，列号即出错位置——前面算出的
$S=(0,1,0,1)^{\mathsf T}$ 正是 $H$ 的第 10 列，所以错在第 10 位。

\begin{idea}{教材 5.6.2 与我们这题的关系}
教材用 $r=3$ 的 $(7,4)$ 汉明码，$H$ 是 $3\times7$，每列是 3 位地址，思路完全一样；
而且它是"填满"的（$2^{3}-1=7$，所有非零 3 位地址都用上），故称\textbf{完备码}。
我们这题 $k=8$ 得 $r=4$、$(12,8)$，$2^{4}=16>13$，有 13、14、15 三个地址空着，
所以不是完备码。\textbf{构造方法一字不差，只是矩阵大小不同}。
\end{idea}

\subsection{手把手：已知转移矩阵，分别用两种准则译码（作业 5.11 类型）}

这题与汉明码无关，但同属第 5 章作业且讲评点名过，一并处理。\textbf{题目}：
离散无记忆信道，转移矩阵（行是输入 $x$，列是输出 $y$）为
\[
\bigl(q(y|x)\bigr)=
\begin{pmatrix}
\dfrac{1}{2}&\dfrac{1}{3}&\dfrac{1}{6}\\[2mm]
\dfrac{1}{6}&\dfrac{1}{2}&\dfrac{1}{3}\\[2mm]
\dfrac{1}{3}&\dfrac{1}{6}&\dfrac{1}{2}
\end{pmatrix},
\qquad
p(x)=\bigl(0.5,\ 0.25,\ 0.25\bigr).
\]
求最佳（最大后验）译码规则、极大似然译码规则，以及各自的平均差错率。

\textbf{先把矩阵读顺}：第 1 行第 2 列那个 $\frac13$ 意为"发 $x_1$ 而收到 $y_2$ 的概率是
$\frac13$"，记作 $q(y_2|x_1)=\frac13$。每行加起来都是 1（$\frac12+\frac13+\frac16=1$），
符合"发出一个符号必定收到某个符号"。注意它\textbf{不对称}：$q(y_2|x_1)=\frac13$ 而
$q(y_1|x_2)=\frac16$，所以不能偷懒套对称信道的结论。

\subsubsection{极大似然译码：只比 $q(y|x)$，每列取最大}

\textbf{规则}：收到 $y$ 后把这一列三个数比一比，\textbf{挑最大的}，它对应的 $x$ 就是判决。
注意：\textbf{只看 $q(y|x)$，不许乘 $p(x)$}（讲评原话）。

\begin{center}\small
\setlength{\tabcolsep}{6pt}
\begin{tabular}{@{}lcccc@{}}
\toprule
收到 $y$ & $q(y|x_1)$ & $q(y|x_2)$ & $q(y|x_3)$ & 挑最大 $\Rightarrow$ 判成\\
\midrule
$y_1$ & $\mathbf{1/2}$ & $1/6$ & $1/3$ & $x_1$\\
$y_2$ & $1/3$ & $\mathbf{1/2}$ & $1/6$ & $x_2$\\
$y_3$ & $1/6$ & $1/3$ & $\mathbf{1/2}$ & $x_3$\\
\bottomrule
\end{tabular}
\end{center}

每列最大值恰好都在对角线上（都是 $\frac12$），故 $g_{\mathrm{ML}}(y_1)=x_1$，
$g_{\mathrm{ML}}(y_2)=x_2$，$g_{\mathrm{ML}}(y_3)=x_3$，即"收到什么就判成什么"。

\textbf{平均差错率}用"1 减译对的总概率"最好算：
\[
P_{e,\mathrm{ML}}=1-\sum_{y}p\bigl(g(y)\bigr)\,q\bigl(y\,|\,g(y)\bigr).
\]
人话：\textbf{每个收到值，用真正那个 $x$ 的概率 $p(x)$ 乘"从它发出却收到 $y$"的概率，
加起来是译对的概率}，1 减掉就是错的概率。逐项：
\begin{align*}
y_1:\ & p(x_1)q(y_1|x_1)=\tfrac12\times\tfrac12=\tfrac14=0.25,\\
y_2:\ & p(x_2)q(y_2|x_2)=\tfrac14\times\tfrac12=\tfrac18=0.125,\\
y_3:\ & p(x_3)q(y_3|x_3)=\tfrac14\times\tfrac12=\tfrac18=0.125 .
\end{align*}
译对总概率 $=0.25+0.125+0.125=0.5$，故 $P_{e,\mathrm{ML}}=1-0.5=\boxed{0.5}$。

\textbf{换个算法验一遍}（按"发的是谁"分）：$x_1$ 只在收到 $y_2,y_3$ 时判错，
$\frac12\times(\frac13+\frac16)=\frac14$；$x_2$ 在收到 $y_1,y_3$ 时判错，
$\frac14\times(\frac16+\frac13)=\frac18$；$x_3$ 在收到 $y_1,y_2$ 时判错，
$\frac14\times(\frac13+\frac16)=\frac18$。相加 $\frac14+\frac18+\frac18=\frac12$，一致。

\subsubsection{最大后验（最佳）译码：比 $p(x)\,q(y|x)$}

\textbf{规则}：收到 $y$ 后比的是\textbf{联合量} $p(x)q(y|x)$（等价地比 $p(x|y)$，
因为 $p(x|y)$ 只是 $p(x)q(y|x)$ 除以 $p(y)$，而 $p(y)$ 对同一次收到的 $y$ 是常数，
不影响谁最大），挑最大的那个 $x$。

\begin{center}\small
\setlength{\tabcolsep}{5pt}
\begin{tabular}{@{}lcccc@{}}
\toprule
收到 $y$ & $p(x_1)q(y|x_1)$ & $p(x_2)q(y|x_2)$ & $p(x_3)q(y|x_3)$ & 判成\\
\midrule
$y_1$ & $\tfrac12\cdot\tfrac12=\tfrac14$ & $\tfrac14\cdot\tfrac16=\tfrac1{24}$ & $\tfrac14\cdot\tfrac13=\tfrac1{12}$ & $\mathbf{x_1}$\\
$y_2$ & $\tfrac12\cdot\tfrac13=\mathbf{\tfrac16}$ & $\tfrac14\cdot\tfrac12=\tfrac18$ & $\tfrac14\cdot\tfrac16=\tfrac1{24}$ & $\mathbf{x_1}$\\
$y_3$ & $\tfrac12\cdot\tfrac16=\tfrac1{12}$ & $\tfrac14\cdot\tfrac13=\tfrac1{12}$ & $\tfrac14\cdot\tfrac12=\mathbf{\tfrac18}$ & $\mathbf{x_3}$\\
\bottomrule
\end{tabular}
\end{center}

故 $g_{\mathrm{MAP}}(y_1)=x_1$，$g_{\mathrm{MAP}}(y_2)=x_1$，$g_{\mathrm{MAP}}(y_3)=x_3$。

\textbf{与极大似然只差一处：$y_2$ 的判决变了。} 原来 $q(y_2|x_2)=\frac12$ 是列最大，
但乘上 $p(x_2)=0.25$ 后只剩 $\frac18$；而 $x_1$ 虽然 $q(y_2|x_1)=\frac13$ 不是列最大，
可它的 $p(x_1)=0.5$ 最大，乘出来 $\frac16>\frac18$ 于是反超。
\textbf{这就是"先验概率把结论掰过来"的活例子}，也是最大后验优于极大似然的原因：
它知道 $x_1$ 本来就更容易出现。

\textbf{平均差错率}同样用 1 减译对概率：
\begin{align*}
y_1:\ & p(x_1)q(y_1|x_1)=\tfrac14,\\
y_2:\ & p(x_1)q(y_2|x_1)=\tfrac12\times\tfrac13=\tfrac16,\\
y_3:\ & p(x_3)q(y_3|x_3)=\tfrac14\times\tfrac12=\tfrac18 .
\end{align*}
译对总概率 $=\frac14+\frac16+\frac18=\frac{6}{24}+\frac{4}{24}+\frac{3}{24}=\frac{13}{24}$，故
\[
P_{e,\mathrm{MAP}}=1-\frac{13}{24}=\boxed{\frac{11}{24}\approx0.4583}.
\]
验算（按发送方分类）：$x_1$ 只有收到 $y_3$ 时判错，$\frac12\times\frac16=\frac{2}{24}$；
$x_2$ 三个 $y$ 全判成别人，$\frac14\times1=\frac{6}{24}$；$x_3$ 在收到 $y_1,y_2$ 时判错，
$\frac14\times(\frac13+\frac16)=\frac18=\frac{3}{24}$。相加 $\frac{2+6+3}{24}=\frac{11}{24}$，一致。

\begin{exam}{讲评点名的那句话：极大似然别乘 $p(x)$}
这道题最常见的错法是算极大似然时悄悄乘上了 $p(x)$，于是把 $P_{e,\mathrm{ML}}$ 算成 $11/24$
——\textbf{那是最大后验的答案}。区别只有一句话：\textbf{极大似然只看 $q(y|x)$、每列挑最大，
默认各 $x$ 等可能出现}；\textbf{最大后验看 $p(x)q(y|x)$，把"谁更容易出现"也算进去}。
所以\textbf{当且仅当 $p(x)$ 等概时两者给出同一套规则}。本题 $p(x)=(0.5,0.25,0.25)$ 不等概，
答案自然不同：$0.5$ 与 $11/24$，后者更小，符合"最佳译码不差于极大似然"。
\end{exam}

\subsection{考试怎么考}

\begin{exam}{构造汉明码的完整答题步骤（照抄这四个小标题，分步给分）}
\textbf{1. 定 $r$ 与 $n$。} 写不等式 $2^{r}\ge n+1$ 且 $n=k+r$，然后\textbf{把 $r=1,2,3,\dots$
逐个代进去试}，每步写清"此时 $n=?$，$2^{r}=?$ 与 $n+1=?$ 谁大"，最后写"故取 $r=\square$，
$n=k+r=\square$"。\emph{只写答案不写试算过程会扣分。}

\textbf{2. 列位置表。} 列表列出 $1\sim n$ 位，标清哪几位是 $C_1,C_2,C_4,C_8$（2 的幂次位）、
哪几位是 $D_1\sim D_k$（其余按顺序）。\emph{这张表是后面所有计算的依据，必须写全。}

\textbf{3. 写校验方程。} 把每个位置号写成 $r$ 位二进制地址，按"地址第 $j$ 位为 1"分成 $r$ 组，
写 $r$ 个方程（偶校验 $=0$、奇校验 $=1$），再解成 $C_{2^{j}}=\bigoplus(\text{组内信息位})$。

\textbf{4. 代入求校验位、拼码字、回代验算。} 逐个算 $C$（写出异或中间步骤），按位置表拼成
$n$ 位，\textbf{最后一定回代 $r$ 个方程、确认每组 1 的个数满足奇偶要求}，并写一句"验算通过"。

\medskip
\textbf{若还要纠错}：多写一步——把收到的 $y$ 代进同样的 $r$ 个方程得综合征
$(S_{1},S_{2},S_{4},S_{8})$，按权值 $1,2,4,8$ 组装成十进制位置号，指出第几位错并翻转。
\end{exam}

\begin{exam}{四个常见扣分点}
\textbf{1. 校验位位置放错。} 必须放在第 $1,2,4,8,\dots$ 位。若把校验位全堆在末尾
（前 8 位数据、后 4 位校验），那就不是教材这套汉明码了，位置表、分组、地址全部对不上。
\emph{对策：动手算之前先把位置表写出来给阅卷老师看。}

\textbf{2. 偶校验 / 奇校验混淆。} 偶校验要求"每组 1 的个数为偶数"、方程写 $=0$；
奇校验要求"奇数"、方程写 $=1$。\emph{混了就有一半校验位全反，白算。}

\textbf{3. 回代不验算。} 只写最终码字、不代回方程检查，一个位置填错就整题错。
\emph{对策：算完逐组数 1 的个数，把"4/4/2/2，全偶，通过"这行写上去}——既是检查也是给分点。

\textbf{4. 把 $2^{r}\ge n+1$ 写错。} 常见三种：写成 $2^{r}\ge n$（少 1，会漏掉"没错"这个状态）；
忘了 $n$ 里含 $r$（应 $n=k+r$）；不等号方向写反。\emph{对策：写不等式时旁边注明
"$n+1$ 是因为要区分'没错'+$n$ 个'第几位错'"。}
\end{exam}

\subsection{你可能想问}

\begin{qa}{为什么校验位一定要放在 1、2、4、8 位？放在别处不行吗？}
不是"不行"，而是"放这里最省事、最漂亮"。\textbf{第一，只有放在 2 的幂次位，位置的二进制
地址才刚好各缺一位。} 第 1 位地址 $0001$、第 2 位 $0010$、第 4 位 $0100$、第 8 位 $1000$，
它们各自\textbf{只属于一个组}，于是 $C_1$ 只负责 $G_1$ 一组的奇偶性，责任划分干净。
若把校验位放到第 3 位（地址 $0011$），它同时属于 $G_1$ 和 $G_2$，一错就同时弄脏两组结果，
综合征不再"干净地等于出错位置"，纠错就乱了。

\textbf{第二，唯一性。} 要求是"$n$ 个位置各有一个不同的非零地址，好当门牌号"；
$1,2,4,8,\dots$ 正好是能自由取值、互不冲突的位置，而数据位塞在剩下的位置就自动获得
"被多个组同时监控"的待遇，一错有多组报警，地址才排得出来。所以这不是规定，
而是\textbf{让"综合征 = 出错位置"这个性质成立的必要条件}。
\end{qa}

\begin{qa}{偶校验和奇校验有什么区别？选哪个？}
\textbf{区别只一处}：偶校验要求每组 1 的个数是\textbf{偶数}（方程右端 $=0$），奇校验要求是
\textbf{奇数}（方程右端 $=1$）；于是求校验位时"补差"方向反过来——组内信息位已有奇数个 1 时，
偶校验要 $C=1$，奇校验要 $C=0$。

\textbf{纠错能力完全一样}，都能纠 1 位错、$d_{\min}$ 都是 3。区别只在"全 0 码字"：
偶校验下 \texttt{00}$\cdots$\texttt{0} 是合法码字（0 个 1 是偶数），奇校验下不是，
所以奇校验能顺带发现"全 0"这种明显异常（比如线路断了收不到信号）。工程上两种都常见，
\textbf{考试时题目说哪个就用哪个}，把方程右端写成对应的 0 或 1 即可，其余步骤一字不改。
\end{qa}

\begin{qa}{汉明码既然厉害，为什么只纠 1 位错？错 2 位怎么办？}
因为它 $d_{\min}=3$，由 $d_{\min}\ge2t+1$ 得 $t=1$，\textbf{这是距离决定的硬上限，
跟算法聪不聪明无关}。错 2 位时两个错误会弄脏两组校验结果，综合征指向"第 3 个位置"——
既不是真正的出错位置也不表示无错，收方会\textbf{自信地纠错，结果越纠越错}。
这是汉明码的固有风险：\textbf{它能纠 1 位错，但错 2 位时可能纠歪。}

要纠更多错就得拉开距离：把汉明码再加一位总校验（扩展汉明码）可把 $d_{\min}$ 从 3 提到 4，
从而"纠 1 位 + 检出 2 位"；要真纠 2 位需 $d_{\min}\ge5$，得用 BCH 码、RS 码、卷积码这类
更强的码，代价是码率进一步下降——又回到那句老话：\textbf{可靠性是用长度换的}。
\end{qa}

\begin{idea}{这一讲的三句话总结}
\textbf{1. 为什么要编码}：线会翻错 bit，裸传原码收方无从察觉；加冗余位（校验位）让错误留下
痕迹。冗余 = 用更长的码换可靠，方向与信源编码正好相反。

\textbf{2. 汉明码怎么造}：$2^{r}\ge n+1$ 定 $r$ $\to$ 校验位放 $1,2,4,8$ 位 $\to$ 按二进制
地址分 $r$ 组写偶校验方程 $\to$ 代入求校验位、回代验算
（$k=8$ 得 $r=4,n=12$；\texttt{10011100} $\to$ \texttt{111100101100}）。

\textbf{3. 怎么纠错}：把收到的 $y$ 代进同样 $r$ 个方程得综合征，\textbf{综合征就是出错位置的
二进制地址}，翻回来即可（$(0,1,0,1)\to$ 第 10 位）。
\end{idea}
